
\section{Aufgabe 1}

Die Messungen wurden auf einer NVIDIA GeForce RTX 4060 durchgeführt. Die
verwendete GPU besitzt 24 SMs, 3072 FP32-Kerne, einen FP32-Peak von
ungefähr 15{,}3 TFLOP/s und eine Speicherbandbreite von 272 GB/s.

\subsection{Verwendete Kernel}

\subsubsection{FMA-Kernel für Skalierung und GPU-Warm-up}

Der erste verwendete Kernel ist ein \textit{fused multiply add (fma)} Kernel: 
\texttt{kernel\_fma} in \texttt{fma.cu}. Der Kernel simuliert eine große arithmetische 
Last bei gleichzeitig wenigen Speicherzugriffen. Damit kann gemessen werden, 
ab welcher Problemgröße genügend Threads vorhanden sind, um die FP32-Recheneinheiten 
der GPU voll auszulasten.\par

Jeder Thread bearbeitet genau eine Element, d.i. er lädt einen 4 Byte großen Wert, 
führt mindestens eine FMA Operation auf diesem Wert aus und schreibt anschließend das Ergebnis zurück. 
Dadurch ist der Speicherbedarf minimal (4 Byte \textit{read} und 4 Byte \textit{write}), 
während die arithmetische Last über die Anzahl der Iterationen der FMA Operation mithilfe 
des Parameters \texttt{iters} skaliert werden kann.

\begin{verbatim}
float x = in[idx];
float acc0 = x, acc1 = x * 0.5f + 1.0f;

for (int i = 0; i < iters; ++i) {
    acc0 = fmaf(acc0, a, b);
    acc1 = fmaf(acc1, a, b);
}

out[idx] = acc0 + acc1;
\end{verbatim}

Das Ergebnis wird in den Ausgabepuffer (\texttt{out[idx]}) geschrieben, damit der Compiler die
Rechenschleife nicht entfernen kann.

\subsubsection{FMA Divergenz-Kernel für unterschiedliche Kontrollflusspfade}

Der zweite verwendete Kernel ist ebenfalls ein \textit{fused multiply add (fma)} Kernel: 
\texttt{kernel\_fma\_div} in \texttt{fma.cu}. 
Im Unterschied zum 1. Kernel fügt dieser gezielt einen datenabhängigen Kontrollfluss hinzu. 
Damit kann nicht nur die reine Rechenleistung gemessen werden, sondern auch der Effekt, der entsteht, 
wenn verschiedene \textit{Lanes} eines Warps unterschiedliche Pfade nehmen können.

Im Unterschied zu \texttt{kernel\_fma} ist die Arbeit pro Thread in \texttt{kernel\_fma\_div} 
nicht uniform verteilt. Einige Threads nehmen einen schweren Pfad mit vielen FMA-Operationen,
während andere einen leichten Pfad mit deutlich weniger Arbeit nehmen. Der Parameter
\texttt{stride} steuert dabei, ob diese Pfade innerhalb eines Warps gemischt sind
oder ob ganze Warps homogen bleiben.

\subsubsection{Mandelbrot Kernel der das natürliche Auftreten von Divergenz simuliert}

Der Kernel in \texttt{mandelbrot.cu} berechnet für jeden Pixel eine
Mandelbrot-Iteration

\begin{verbatim}
z = z * z + c
\end{verbatim}

bis entweder \texttt{|z| > 2} gilt oder \texttt{maxIter} erreicht ist. Es handelt sich 
hierbei ebenfalls um ein rechenintensives Problem mit sehr wenigen Speicherzugriffen, da
jeder Thread am Ende im Wesentlichen nur die erreichte Iterationszahl schreibt.

Der Unterschied zum FMA-Kernel ist, dass die Divergenz
nicht künstlich über \texttt{stride} eingestellt wird. Sie entsteht
natürlich aus den Daten: Pixel außerhalb der Mandelbrot-Menge brechen früh
ab, Pixel nahe am Rand laufen deutlich länger, und Pixel im Inneren erreichen
oft \texttt{maxIter}. Wenn solche Pixel in derselben Warp liegen, müssen die
unterschiedlich langen Schleifenpfade serialisiert werden. Mandelbrot zeigt
also denselben SIMT-Effekt in einem realistischeren Algorithmus als der FMA-Kernel, ist aber
weniger exakt kontrollierbar als \texttt{kernel\_fma\_div}.

\subsubsection{Monte Carlo Kernel um das zufällige Auftreten von Divergenz zu simulieren}

Der Monte-Carlo Kernel in \texttt{montecarlo.cu} schätzt \(\pi\), indem jeder Thread 
zufällige Punkte im Einheitsquadrat erzeugt und prüft, ob sie im Viertelkreis liegen:

\begin{verbatim}
if (x * x + y * y <= 1.0f)
    ++local;
\end{verbatim}

Dabei ist die Divergenz nicht räumlich strukturiert wie bei der Berechnung der Mandelbrot Menge, 
sondern sie entsteht zufällig durch die Stichproben. In einer Warp nehmen die Lanes daher 
statistisch gemischte Pfade. 

Monte Carlo ist allerdings kein isolierter Divergenz-Benchmark, da Die Laufzeit
zusätzlich Kosten für die Erzeugung der Zufallszahlen, lokale Zähler und
eine globale Reduktion per \texttt{atomicAdd} enthält. Deshalb dient dieser Kernel nur
als ergänzender Referenzfall, um zu illustrieren, wie zufällige Kontrollflüsse in
einem stochastischen GPU-Algorithmus auftreten.

\subsection{Skalierung über die Problemgröße}

Für die Skalierung wurde die Anzahl der Iterationen fest gewählt: \texttt{iters = 512}. Die
Problemgröße \textit{n} wurde von $1024$ bis $134$ Millionen Elementen
variiert. Die Ergebnisse liegen in \texttt{outputs/scaling.csv}; die daraus erzeugte
Grafik ist in Abbildung~\ref{fig:scaling} dargestellt.

\begin{figure}[H]
  \centering
  \includegraphics[width=0.95\linewidth]{scaling.png}
  \caption{Durchsatz des FMA-Kernels in Abhängigkeit von der Problemgröße.}
  \label{fig:scaling}
\end{figure}

\begin{table}[H]
  \centering
  \begin{tabular}{>{\raggedleft\arraybackslash}p{3.0cm}
                  >{\raggedleft\arraybackslash}p{3.0cm}
                  >{\raggedleft\arraybackslash}p{4.0cm}}
    \hline
    \textbf{n = Threads} & \textbf{GFLOP/s} & \textbf{Vielfaches der Sättigung} \\
    \hline
    1\,024       & 819    & 0{,}03 \\
    16\,384      & 7\,740 & 0{,}44 \\
    32\,768      & 9\,498 & 0{,}89 \\
    65\,536      & 11\,498 & 1{,}78 \\
    131\,072     & 12\,327 & 3{,}6 \\
    262\,144     & 13\,064 & 7{,}1 \\
    1\,048\,576  & 13\,515 & 28 \\
    16\,777\,216 & 12\,947 & 455 \\
    67\,108\,864 & 12\,536 & 1820 \\
    \hline
  \end{tabular}
  \caption{Auszug aus der Skalierungsmessung.}
  \label{tab:scaling}
\end{table}

Die theoretische Sättigungsschwelle der verwendeten GPU liegt somit bei
\[
24\ \mathrm{SMs} \cdot 1536\ \mathrm{Threads/SM} = 36\,864
\]
gleichzeitig laufenden Threads. Diese Schwelle ist jedoch nicht identisch mit
dem praktischen Plateau. In der Messung wurde die 90-Prozent-Region erst bei
etwa 131\,072 Threads erreicht, also bei ungefähr dem
3{,}6-fachen der Thread-Kapazität. Das Plateau liegt bei etwa
13{,}6 TFLOP/s.

Die Messung zeigt damit den typischen GPU-Warm-up: Bei kleinen
Problemgrößen ist die GPU nicht ausgelastet. Erst wenn die Problemgröße groß genug ist, 
d.i. ausreichend viele Warps und genügend Arbeit vorhanden ist, kann der Scheduler Latenzen durch
Wechseln auf andere lauffähige Warps ausgleichen. Dagegen sind die Kerne einer CPU
auf niedrige Latenz, Sprungvorhersage, Out-of-Order-Ausführung und
große Caches optimiert, sodaß ein einzelner Thread bereits sinnvoll Arbeit leisten kann.
Im Unterschied zur CPU muss die GPU mit sehr vielen unabhängigen Threads versorgt werden.

\subsection{Gezielte Warp-Divergenz}

Auch für den Divergenzversuch mit \texttt{kernel\_fma\_div} hat ein Thread genau ein Element bearbeitet,
wobei die Rechenlast dieses Mal von der Thread-ID abhing. Dabei hat die Hälfte der Threads 
einen schweren Pfad mit \texttt{iters} FMA Operationen, die andere Hälfte
einen leichten Pfad mit \texttt{iters/8} FMA Operationen genommen. Der Parameter \texttt{stride}
hat festgelegt, wie diese Pfade auf die Lanes eines Warps verteilt wurden.

\begin{verbatim}
bool heavy = ((idx / stride) & 1) == 0;

if (heavy) {
    for (int i = 0; i < iters; ++i)
        acc = fmaf(acc, a, b);
} else {
    for (int i = 0; i < iters / 8; ++i)
        acc = fmaf(acc, a, b);
}
\end{verbatim}

Während sich bei \texttt{stride = 1} schwere und leichte Threads innerhalb eines
Warps maximal stark abwechseln, sind Warps bei \textit{strides}, die Vielfache
von 32 sind, homogen. Strides knapp oberhalb von 32 schneiden dagegen wieder
durch Warp-Grenzen und erzeugen teilweise divergente Warps. Dadurch kann direkt
gemessen werden, ob der Performance-Einbruch wirklich ein Intra-Warp-Effekt ist.

Die Messdaten liegen in \texttt{outputs/divergence.csv}; die daraus erzeugte Grafik
ist in Abbildung~\ref{fig:divergence} dargestellt.

\begin{figure}[H]
  \centering
  \includegraphics[width=0.95\linewidth]{divergence.png}
  \caption{Performance-Einbruch durch Warp-Divergenz.}
  \label{fig:divergence}
\end{figure}

\begin{table}[H]
  \centering
  \begin{tabular}{>{\raggedleft\arraybackslash}p{2.0cm}
                  >{\raggedleft\arraybackslash}p{3.0cm}
                  >{\raggedleft\arraybackslash}p{3.0cm}
                  p{3.0cm}}
    \hline
    \textbf{Stride} & \textbf{GFLOP/s} & \textbf{Zeit [ms]} & \textbf{Regime} \\
    \hline
    1  & 5\,778  & 13{,}38 & max. divergent \\
    2  & 5\,822  & 13{,}28 & divergent \\
    4  & 5\,913  & 13{,}08 & divergent \\
    8  & 5\,739  & 13{,}47 & divergent \\
    16 & 6\,016  & 12{,}85 & divergent \\
    32 & 11\,264 & 6{,}86  & divergenzfrei \\
    33 & 6\,170  & 12{,}53 & teils divergent \\
    34 & 6\,384  & 12{,}11 & teils divergent \\
    36 & 6\,620  & 11{,}68 & teils divergent \\
    40 & 7\,334  & 10{,}54 & teils divergent \\
    48 & 8\,525  & 9{,}07  & teils divergent \\
    64 & 11\,065 & 6{,}99  & divergenzfrei \\
    65 & 7\,511  & 10{,}29 & teils divergent \\
    66 & 7\,732  & 10{,}00 & teils divergent \\
    68 & 7\,901  & 9{,}78  & teils divergent \\
    72 & 8\,315  & 9{,}30  & teils divergent \\
    80 & 9\,039  & 8{,}55  & teils divergent \\
    96 & 10\,684 & 7{,}24  & divergenzfrei \\
    \hline
  \end{tabular}
  \caption{Divergenzmessung bei gleicher Gesamtarbeit.}
  \label{tab:divergence}
\end{table}

Die Auswertung zeigt, dass der Durchsatz genau bei \textit{strides} springt, die
Vielfache der Warp-Größe $32$ sind. Für \texttt{stride < 32} enthält jeder Warp
beide Kontrollflusspfade, sodaß der Warp diese Pfade serialisieren muss. Der
gemessene Durchsatz liegt in diesem Fall nur bei etwa 5{,}7 bis
6{,}0 TFLOP/s. Bei \texttt{stride = 32}, \texttt{64} und \texttt{96} sind die
Warps homogen; der Durchsatz steigt auf etwa 10{,}7 bis 11{,}3 TFLOP/s.
Die Zwischenwerte \texttt{33} bis \texttt{48} und \texttt{65} bis
\texttt{80} schneiden die Warp-Grenzen nicht sauber und sind deshalb nur
teilweise divergent. Die vollständig divergente Variante ist gegenüber den
warp-homogenen Fällen ungefähr um den Faktor 1{,}9 langsamer.

\subsection{Warum Divergenz auf der GPU teuer ist}

Eine GPU führt Threads in Warps aus, die in der getesten GPU aus 32 Lanes bestehen. 
Im SIMT-Modell müssen alle Lanes gemeinsam einen Instruktionsstrom ausführen. Das heißt, wenn einzelne
Lanes unterschiedliche Pfade eines \texttt{if}/\texttt{else} nehmen, kann die
GPU diese Pfade nicht gleichzeitig als unabhängige Kontrollflüsse
ausführen. Stattdessen wird zuerst einer der beiden Pfade mit einer Aktivitätsmaske
ausgeführt, danach der andere Pfad mit der umgekehrten Maske. Lanes, die zum
jeweils anderen Pfad gehören, sind in dieser Zeit inaktiv.

Dadurch addieren sich die Laufzeiten der Pfade innerhalb des Warps. Im
Extremfall muss ein Warp viele verschiedene Pfade nacheinander ausführen. 
Genau dieser Effekt ist in Abbildung~\ref{fig:divergence} sichtbar:
Der Performance-Sprung liegt nicht bei beliebigen Blockgrößen, sondern
exakt an den Vielfachen der Warp-Grenze von 32 Threads.

\subsection{Warum dieselbe Verzweigung auf der CPU weniger stört}

Eine CPU besitzt kein Warp-Modell mit 32 aneinander gekoppelten Lanes. Stattdessen 
führt ein CPU-Kern einen einzelnen Instruktionsstrom aus und reagiert auf Verzweigungen
mit Sprungvorhersage und spekulativer Ausführung. Bei gut
vorhersagbaren Mustern kostet eine Verzweigung daher fast nichts. Zudem besitzt jeder
CPU-Kern jeweils einen eigenen Kontrollfluss.

Der Architekturunterschied ist deshalb grundsätzlich: Die GPU spart
Steuerlogik und erzielt bei einem uniformen Kontrollfluss einen sehr hohen Durchsatz, 
muss dafür aber bei einem divergenten Kontrollfluss innerhalb eines Warps serialisieren. 
Dagegen investiert die CPU wesentlich mehr Logik pro Kern, um auch einen komplexen 
Kontrollfluss mit niedriger Latenz auszuführen, kann dafür allerdings nicht den Durchsatz 
einer GPU bei einem uniformen Kontrollflusserreichen.

\subsection{Anhang: Ergänzende Referenzkernel}

Zusätzlich zu den synthetischen FMA- und Divergenz-Kerneln wurden der
Mandelbrot- und der Monte-Carlo-Kernel als ergänzende Referenzfälle ausgewertet.

\subsubsection{Mandelbrot-Kernel}

Der Mandelbrot-Kernel berechnet für jeden Pixel, nach wie vielen Iterationen
die Folge $z_{n+1}=z_n^2+c$ den Radius $2$ überschreitet. Dadurch hat jeder
Thread eine datenabhängige Laufzeit. Punkte außerhalb der Menge brechen früh
ab, Punkte nahe am Rand oder innerhalb der Menge laufen deutlich länger.

\begin{figure}[H]
  \centering
  \includegraphics[width=0.92\linewidth]{mandelbrot_scaling.png}
  \caption{Mandelbrot-Sweep: Laufzeit und mittlere Iterationszahl pro Pixel.}
  \label{fig:mandelbrot_scaling}
\end{figure}

\begin{table}[H]
  \centering
  \begin{tabular}{>{\raggedleft\arraybackslash}p{2.4cm}
                  >{\raggedleft\arraybackslash}p{2.4cm}
                  >{\raggedleft\arraybackslash}p{2.2cm}
                  >{\raggedleft\arraybackslash}p{2.4cm}
                  >{\raggedleft\arraybackslash}p{2.4cm}}
    \hline
    \textbf{Auflösung} & \textbf{Pixel} & \textbf{Zeit [ms]} &
    \textbf{GFLOP/s, obere Schranke} & \textbf{mittl. Iter.} \\
    \hline
    $512^2$  & 262\,144    & 0{,}114 & 22\,939 & 177{,}6 \\
    $1024^2$ & 1\,048\,576 & 0{,}361 & 29\,008 & 177{,}5 \\
    $2048^2$ & 4\,194\,304 & 1{,}330 & 31\,527 & 177{,}4 \\
    $4096^2$ & 16\,777\,216 & 5{,}115 & 32\,797 & 177{,}4 \\
    \hline
  \end{tabular}
  \caption{Mandelbrot-Messung mit \texttt{maxIter = 1000}. Die GFLOP/s-Spalte
  ist eine obere Schranke, da sie für jedes Pixel die maximale Iterationszahl
  annimmt.}
  \label{tab:mandelbrot}
\end{table}

\begin{figure}[H]
  \centering
  \includegraphics[width=0.88\linewidth]{mandelbrot_result.png}
  \caption{Mandelbrot: Die Farbe kodiert die Anzahl der Iterationen bis zum
  Abbruch.}
  \label{fig:mandelbrot_result}
\end{figure}

Die Messung zeigt, dass die mittlere Iterationszahl über die skalierten
Auflösungen nahezu konstant bleibt, während die Laufzeit mit der Pixelzahl
wächst. Das Bild illustriert die Ursache der Divergenz sichtbar:
benachbarte Pixel können sehr unterschiedliche Iterationszahlen besitzen,
insbesondere am Rand der Menge.

\subsubsection{Monte-Carlo-Kernel}

Der Monte-Carlo-Kernel schätzt $\pi$, indem zufällige Punkte im Einheitsquadrat
geraten werden. Ein Punkt zählt als Treffer, wenn $x^2+y^2 \leq 1$ gilt. Die
Divergenz hängt in diesem Fall direkt von den zufällig gezogenen Zahlen ab.

\begin{figure}[H]
  \centering
  \includegraphics[width=0.92\linewidth]{montecarlo_scaling.png}
  \caption{Monte-Carlo-Sweep: Stichprobendurchsatz und Fehler der
  $\pi$-Schätzung.}
  \label{fig:montecarlo_scaling}
\end{figure}

\begin{table}[H]
  \centering
  \small
  \setlength{\tabcolsep}{4pt}
  \begin{tabular}{>{\raggedleft\arraybackslash}p{2.4cm}
                  >{\raggedleft\arraybackslash}p{2.6cm}
                  >{\raggedleft\arraybackslash}p{2.0cm}
                  >{\raggedleft\arraybackslash}p{2.2cm}
                  >{\raggedleft\arraybackslash}p{2.2cm}
                  >{\raggedleft\arraybackslash}p{2.2cm}}
    \hline
    \textbf{Samp./Th.} & \textbf{Samples [Mrd.]} &
    \textbf{Zeit [ms]} & \textbf{Gsamp./s} & \textbf{$\hat{\pi}$} &
    \textbf{Fehler} \\
    \hline
    256  & 0{,}268 & 1{,}969  & 136{,}36 & 3{,}141646 & $5{,}36\cdot10^{-5}$ \\
    512  & 0{,}537 & 3{,}913  & 137{,}20 & 3{,}141678 & $8{,}52\cdot10^{-5}$ \\
    1024 & 1{,}074 & 7{,}679  & 139{,}82 & 3{,}141688 & $9{,}49\cdot10^{-5}$ \\
    2048 & 2{,}147 & 14{,}172 & 151{,}53 & 3{,}141680 & $8{,}76\cdot10^{-5}$ \\
    4096 & 4{,}295 & 28{,}391 & 151{,}28 & 3{,}141648 & $5{,}57\cdot10^{-5}$ \\
    8192 & 8{,}590 & 56{,}752 & 151{,}36 & 3{,}141613 & $2{,}02\cdot10^{-5}$ \\
    \hline
  \end{tabular}
  \caption{Monte-Carlo-Messung mit $2^{20}$ Threads.}
  \label{tab:montecarlo}
\end{table}

\begin{figure}[H]
  \centering
  \includegraphics[width=0.72\linewidth]{montecarlo_result.png}
  \caption{Monte-Carlo-Ergebnisbild: zufällige Punkte im Einheitsquadrat,
  getrennt nach Treffern innerhalb des Viertelkreises und Punkten außerhalb.}
  \label{fig:montecarlo_result}
\end{figure}

Die Tabelle zeigt, dass der Durchsatz ab größeren Stichprobenzahlen pro Thread
bei etwa $151$ Gsamples/s liegt. Der Fehler der $\pi$-Schätzung sinkt nicht strikt
monoton, da die Messung stochastisch ist; mit wachsender Gesamtzahl an Samples
wird die Schätzung aber tendenziell stabiler.

%Reproduzierbar sind die Messungen mit:
%
%\begin{verbatim}
%make
%make all
%make scaling
%make divergence
%make mandelbrot
%make montecarlo
%make report
%\end{verbatim}
